\clearpage
% C058：保留原三道题并完善取等验证；增加方向推广、真假逆命题辨析。
\begin{examplebox}
\textbf{\textcolor{CExample}{例 1（基础：已知一段求另一段）}}

\textbf{题目：}抛物线 $y^2=8x$ 的焦点为 $F$，过 $F$ 的弦 $AB$ 满足 $|AF|=6$，求 $|BF|$。

\textbf{解：}由 $2p=8$，得 $p=4$，从而
\[
 \frac16+\frac1{|BF|}=\frac2p=\frac12
 \quad\Longrightarrow\quad \frac1{|BF|}=\frac13
 \quad\Longrightarrow\quad \boxed{|BF|=3}.
\]
\textbf{为什么如此快：}题干已有“过焦点的弦”，可\textbf{直接用倒数和}，无需先求两个交点。

\medskip
\textbf{\textcolor{CExample}{例 2（逆向求参数：已知两段）}}

\textbf{题目：}抛物线 $y^2=2px$（$p>0$）的焦点弦 $AB$ 满足 $|AF|=3$、$|BF|=6$，求 $p$ 与抛物线方程。

\textbf{解：}
\[
 \frac13+\frac16=\frac2p
 \quad\Longrightarrow\quad \frac12=\frac2p
 \quad\Longrightarrow\quad \boxed{p=4},\qquad\boxed{y^2=8x}.
\]
\textbf{防错：}本题的“逆向”只是由\textbf{已知焦点弦}求参数，并非用倒数和判定未知直线是不是焦点弦。

\medskip
\end{examplebox}
\clearpage
\begin{examplebox}
\textbf{\textcolor{CExample}{例 3（最值：弦长下界与取等）}}

\textbf{题目：}抛物线 $y^2=4x$ 的焦点弦 $AB$，求 $|AB|$ 的最小值及对应直线。

\textbf{解：}由 $2p=4$ 得 $p=2$。令 $r=|AF|,s=|BF|$，由于 $F$ 在 $A,B$ 间，$|AB|=r+s$。又
\[
 \frac1r+\frac1s=\frac2p=1,
 \quad
 (r+s)\left(\frac1r+\frac1s\right)
 =2+\frac rs+\frac sr\ge4,
\]
故 $\boxed{|AB|\ge4}$，等号当且仅当 $r=s$。取等时 $r=s=p=2$，焦半径公式给出 $x_A=x_B=1$；两交点不同且纵坐标相反，因此对应直线
\[
 \boxed{x=1}\quad\text{（过焦点的通径）},\qquad\boxed{|AB|_{\min}=4}.
\]
\textbf{为什么取等必须核查：}均值不等式给出下界后，还要核实确实存在满足 $r=s$ 的焦点弦；通径就是这样的弦。

\medskip
\textbf{\textcolor{CExample}{例 4（改变问法：已知弦长求两段）}}

\textbf{题目：}在 $y^2=8x$ 中，焦点弦 $AB$ 长为 $9$，求两段焦半径（无须区分 $A,B$ 顺序）。

\textbf{解：}由 $p=4$，有 $r+s=9$ 且 $\frac{r+s}{rs}=\frac12$，于是 $rs=18$。因此 $r,s$ 是方程
\[
 z^2-9z+18=0\quad\Longleftrightarrow\quad (z-3)(z-6)=0
\]
的两根，得到 $\boxed{\{r,s\}=\{3,6\}}$。

\textbf{迁移点：}“倒数和”与“弦长是两段的和”合用，可迅速求两段的\textbf{积和根}。

\medskip
\end{examplebox}
\clearpage
\begin{examplebox}
\textbf{\textcolor{CExample}{例 5（方向变式：开口向左也成立）}}

\textbf{题目：}抛物线 $y^2=-12x$ 的焦点弦 $AB$，求 $\frac1{|AF|}+\frac1{|BF|}$。

\textbf{解：}按 $y^2=-2px$（$p>0$）识别 $2p=12$，即 $p=6$。将坐标轴反向即可化为右开口同型式，长度不变，所以
\[
 \boxed{\frac1{|AF|}+\frac1{|BF|}=\frac2p=\frac13}.
\]
\textbf{注意：}开口方向改变，\textbf{不是}把 $p$ 带负号代入；这里的 $p$ 始终取正值。

\medskip
\end{examplebox}
\clearpage
\begin{examplebox}
\begin{center}\begin{tikzpicture}[font=\small,>=stealth]
\draw[gray!45,rounded corners=3pt] (0,0) rectangle (14.2,10.5);
\node at (7.1,9.85) {只把近端点翻到 x 轴上方：距离没变，共线关系变了};
\node at (3.4,8.9) {真正焦点弦：三点共线};
\begin{scope}[shift={(2.0,4.6)},scale=.82]
\draw[teal!75!black,thick,samples=90,domain=-1.4:4.5,variable=\v] plot ({\v*\v/4},{\v});
\coordinate (A) at (.25,-1);\coordinate (B) at (4,4);\coordinate (F) at (1,0);
\draw[gray] (-.5,0)--(5.5,0);
\draw[blue!75!black,thick] (F)--(A);
\draw[orange!85!black,thick] (F)--(B);
\draw[violet,dashed,thick] (A)--(B);
\foreach \q in {A,B,F} \fill (\q) circle(3pt);
\node[left=10pt] at (A) {$A$};\node[above left=6pt] at (B) {$B$};\node[below right=6pt] at (F) {$F$};
\end{scope}
\node[align=center] at (3.4,2.4) {$AF=5/4,\ BF=5$\\[10pt]$\displaystyle\frac1{AF}+\frac1{BF}=1$};
\node at (10.1,8.9) {反例：等式成立，却不共线};
\begin{scope}[shift={(8.7,4.6)},scale=.82]
\draw[teal!75!black,thick,samples=90,domain=-1.4:4.5,variable=\v] plot ({\v*\v/4},{\v});
\coordinate (A) at (.25,1);\coordinate (B) at (4,4);\coordinate (F) at (1,0);
\draw[gray] (-.5,0)--(5.5,0);
\draw[blue!75!black,thick] (F)--(A);
\draw[orange!85!black,thick] (F)--(B);
\draw[violet,dashed,thick] (A)--(B);
\foreach \q in {A,B,F} \fill (\q) circle(3pt);
\node[left=10pt] at (A) {$A$};\node[above left=6pt] at (B) {$B$};\node[below right=6pt] at (F) {$F$};
\end{scope}
\node[align=center] at (10.1,2.4) {$AF=5/4,\ BF=5$\\[10pt]$\displaystyle\frac1{AF}+\frac1{BF}=1$};
\node at (7.1,.65) {右图正是例6：紫虚线 $AB$ 绕开了 $F$；所以不能只凭倒数和反推“过焦点”。};
\end{tikzpicture}
\end{center}
\textbf{\textcolor{CExample}{例 6（反例：倒数和相等不能证明过焦点）}}

\textbf{题目：}在 $y^2=4x$ 中，给出 $A(\frac14,1)$、$B(4,4)$。检验倒数和是否成立，并判断 $AB$ 是否过焦点。

\textbf{解：}抛物线参数 $p=2$，焦点 $F(1,0)$，且 $A,B$ 均在抛物线上；由焦半径公式
\[
 |AF|=\frac14+1=\frac54,\qquad |BF|=4+1=5,
\]
\[
 \frac1{|AF|}+\frac1{|BF|}=\frac45+\frac15=1=\frac2p.
\]
图中紫虚线不经过 $F$。坐标核验：$A,B$ 的直线斜率为 $\frac{4-1}{4-1/4}=\frac45$，代入 $x=1$ 得其直线上 $y=\frac85\ne0$，故\textbf{$F$ 不在直线 $AB$ 上}。

\textbf{结论：}倒数和定值以及 $x_Ax_B=p^2/4$ 是\textbf{焦点弦的必要条件，不是充分条件}。
\end{examplebox}
